\chapter{One Homotopy, Many Names}
\label{ch:homotopy}

Chapter~\ref{ch:parametric} developed the parametric active-set homotopy as a piece
of optimization. This chapter makes the claim that gives Part~III its point: the
homotopy is not an implementation detail of any one method but a \emph{computational
primitive}, and at least four literatures arrived at it independently. In one case
two of the instances are not analogous but \emph{identical}, and
Theorem~\ref{thm:identity} proves it by a one-line substitution.

\section{The instances}

Read through Chapter~\ref{ch:parametric}, each member is the common scheme with its
quadratic form and its face-mechanism filled in.

\paragraph{The LASSO and least angle regression.} $H = X\T X$, $d = X\T y$; the
active set is \emph{induced} by the $\ell_1$ subgradient and is non-monotone. A
coordinate enters when its correlation with the residual reaches $\lambda$ (event
$\mathrm{D}_1$ with a moving threshold) and leaves when its coefficient crosses zero
(event $\mathrm{P}_1$). Forward least angle regression is the entering-only
ancestor; the LASSO modification restores the leave move
\cite{efron2004,osborne2000,tibshirani1996}.

\paragraph{The Critical Line Algorithm.} $H = \Sig$, $d = \mu$; the active set is
\emph{imposed} by the box bounds and is non-monotone, with assets leaving and
re-entering the free set as $\lambda$ decreases, tracing the \emph{corner
portfolios} of the efficient frontier. It is the oldest instance by four decades
\cite{markowitz1956,markowitz1959}.
\index{Critical Line Algorithm}

\paragraph{The support-vector-machine path.} Hastie et al.\ \cite{hastie2004svm}
trace the SVM solution as the cost parameter varies; $H$ is a \emph{kernel} Gram
matrix, the active set is the set of margin support vectors, and the events are
points crossing the margin. Induced and non-monotone, and the sharpest illustration
of the operator reading of Chapter~\ref{ch:operator}: the path touches the kernel
only through products, never forming it densely.

\paragraph{The elastic net and the generalized lasso.} The elastic net
\cite{zou2005} adds a ridge, sending $H \to X\T X + \eta I$ and otherwise leaving the
scheme untouched: the regularisation merely conditions the quadratic form. The
generalized lasso \cite{tibshirani2011}, with penalty $\lambda\norm{D\beta}_1$ for a
structured difference operator---the fused lasso, trend filtering, and graph
penalties as special cases---traces a piecewise-linear \emph{dual} path by the same
homotopy; the face mechanism acts on $D\beta$ rather than $\beta$, a change of
coordinates, not of scheme.

\paragraph{Explicit model predictive control.} The constrained linear-quadratic
regulator solved as a function of the state \cite{bemporad2002,tondel2003} is the
multidimensional cousin of Remark~\ref{rem:vector}: the parameter is a
\emph{vector}, the interval of breakpoints becomes a polyhedral partition into
critical regions, and the homotopy becomes region traversal. Only its one-parameter
slice is literally \eqref{eq:pqp}, so explicit MPC is a \emph{generalisation} of the
common scheme rather than another instance of the scalar identity below.

\medskip
None of these needs machinery beyond Chapter~\ref{ch:parametric}. What looks like
five literatures is one algorithm with five fillings.

\section{The efficient frontier, for the reader who has not met it}
\label{sec:frontier-primer}

Since the identity below pairs finance with regression, a paragraph of each for the
reader who knows only one.

Mean--variance portfolio selection is a bi-objective problem, and the statistician
already knows its shape. An investor spreads weights $w$ over $n$ assets, the budget
$\onevec\T w = 1$ fixing the scale and a long-only mandate adding $w \ge 0$, and
wants two incompatible things: high expected return $\mu\T w$ and low risk
$w\T\Sig w$. No portfolio optimises both, so the object of study is the
\emph{efficient frontier}: the portfolios that are not dominated. That is precisely a
\emph{Pareto front}, and ``Markowitz-efficient'' is Pareto's nondomination in another
vocabulary.
\index{efficient frontier}

A convex Pareto front is recovered by scalarisation: weight the two objectives with a
single $\lambda \ge 0$ and minimise $\tfrac12 w\T\Sig w - \lambda\mu\T w$, which is
\eqref{eq:pqp}. Convex risk and polyhedral constraints make the sweep \emph{exact},
with nothing skipped: the front is an exact parabola in the
(variance, return) plane between corner portfolios, and those corners---where the set
of held assets changes---are its breakpoints.

So the efficient frontier is a Pareto front of the same \emph{kind} as a
regularisation path: a convex bi-objective traded off by one weight and traced corner
to corner. Which front, and traded against what, differs by instance: return against
variance here, fit against sparsity for the LASSO.

\section{The exact identity}
\label{sec:identity}

Now the sharpest special case: two of the most dissimilar members are, under one
substitution, the same instance.

For weights $w \in \R^n$, expected returns $\mu$, and covariance $\Sig \succ 0$, the
gross-exposure-constrained mean--variance program is
\begin{equation}
\label{eq:mark}
\begin{aligned}
\mathrm{M}_c:\quad \min_{w}\ &\tfrac12\, w\trans \Sig\, w - \mu\trans w \\
\text{s.t.}\quad &\norm{w}_1 \le c,\ \ A w = b,\ \ G w \le h .
\end{aligned}
\end{equation}
For a response $y \in \R^m$ and design $X \in \R^{m\times n}$, the constrained LASSO
is
\begin{equation}
\label{eq:lasso}
\begin{aligned}
\mathrm{L}_\lambda:\quad \min_{\beta}\ &\tfrac12\,\norm{y - X\beta}_2^2 + \lambda\,\norm{\beta}_1 \\
\text{s.t.}\quad &A\beta = b,\ \ G\beta \le h .
\end{aligned}
\end{equation}

\begin{theorem}[CLA--LASSO identity under general constraints]
\label{thm:identity}
\index{identity!CLA and LASSO}
Suppose $\Sig = X\T X$ and $\mu = X\T y$. Then for every $\lambda \ge 0$ there is a
$c(\lambda) = \norm{\beta(\lambda)}_1$, nonincreasing in $\lambda$, such that
$\mathrm{M}_{c(\lambda)}$ and $\mathrm{L}_\lambda$ have the same minimiser. The
solution path of $\mathrm{M}_c$ as $c$ varies and the regularization path of
$\mathrm{L}_\lambda$ as $\lambda$ varies are therefore the \emph{same}
piecewise-linear curve in $\R^n$, traced by the same active-set homotopy and sharing
the same breakpoints as points of $\R^n$.
\end{theorem}

\begin{proof}
Completing the square,
\[
  \tfrac12\norm{y-Xw}_2^2 = \tfrac12\, w\T (X\T X)\, w - (X\T y)\T w + \tfrac12\norm{y}_2^2 .
\]
With $\Sig = X\T X$ and $\mu = X\T y$ the objective of $\mathrm{M}_c$ equals
$\tfrac12\norm{y-Xw}_2^2$ up to the additive constant $\tfrac12\norm{y}_2^2$, which
does not affect the minimiser. Hence, under the same linear constraints,
$\mathrm{M}_c$ is exactly the constrained least-squares problem subject to
$\norm{w}_1 \le c$. Its Lagrangian relaxation of that single constraint is
$\mathrm{L}_\lambda$. Because the objective is convex and $\norm{\cdot}_1$ is convex,
strong duality holds for the norm-ball constraint, so the constrained and penalised
forms have identical solution sets as $c$ and $\lambda$ range over $[0,\infty)$, with
$c(\lambda) = \norm{\beta(\lambda)}_1$ the active constraint level at the dual value
$\lambda$; $c$ is nonincreasing because larger penalties shrink $\norm{\beta}_1$.
Both objectives are quadratic and both feasible sets polyhedral, so by
Lemma~\ref{lem:affine} each path is piecewise linear; sharing a minimiser at every
parameter value, the two paths are the same set of points with the same corners, and
tracing either by the homotopy visits the same vertices in the same order.
\end{proof}

\begin{remark}[When the substitution is available]
The hypothesis is mild and runs in both directions. From the regression side it holds
by definition. From the portfolio side, where one starts with $(\Sig,\mu)$, any
$\Sig \succeq 0$ factors as $X\T X$ (a symmetric square root or Cholesky factor
serves), and $\mu = X\T y$ is solvable for $y$ precisely when
$\mu \in \range(X\T)$. When $\Sig \succ 0$---the finance regime---$X$ has full column
rank and a representing $y$ always exists, so the theorem applies to the
mean--variance problem for every $(\Sig,\mu)$. The only obstruction can occur for
singular $\Sig$, and bites in neither instance.
\end{remark}

\subsection{The parameter map}

The map $c(\lambda)$ can be pinned down, which both sharpens the theorem and isolates
the only way the parameter bijection can fail.

\begin{proposition}[Regularity of the $c$--$\lambda$ correspondence]
\label{prop:reg}
Let $\beta(\lambda)$ be the constrained LASSO path and
$c(\lambda) = \norm{\beta(\lambda)}_1$. Then $c$ is continuous and piecewise linear,
affine on each segment with $c'(\lambda) = -s_\Active\T\beta_{\mathrm{slope},\Active}$
over the current support $\Active$ with signs $s_\Active$. In the unconstrained case
$\beta_{\mathrm{slope},\Active} = (X_\Active\T X_\Active)^{-1}s_\Active$, so
$c'(\lambda) < 0$ on every segment with nonempty support and $c$ is a strictly
decreasing, piecewise-linear bijection. In the constrained case
$s_\Active\T\beta_{\mathrm{slope},\Active} \ge 0$, with strict inequality---and hence
the bijection---whenever $s_\Active$ is not annihilated by the reduced Hessian on the
active constraint tangent space.
\end{proposition}

\begin{proof}
On a segment the support and signs are fixed, so
$\norm{\beta(\lambda)}_1 = s_\Active\T(\alpha_\Active - \lambda\beta_{\mathrm{slope},\Active})$,
affine with the stated slope; continuity follows from continuity of $\beta(\lambda)$.
Unconstrained, segment stationarity gives
$\beta_{\mathrm{slope},\Active} = (X_\Active\T X_\Active)^{-1}s_\Active$ and the
quadratic form is positive on a nonempty support. In the constrained case,
eliminating $\zeta$ from~\eqref{eq:border} by its Schur complement expresses the free
block's response to a right-hand side $(g;0)$ as $Z(Z\T H_{\Free\Free}Z)^{-1}Z\T g$,
where the columns of $Z$ span $\nullsp(C_\Free)$. With $g = s_\Active$,
\[
  s_\Active\T\beta_{\mathrm{slope},\Active}
  = \bigl\|(Z\T H_{\Free\Free}Z)^{-1/2}Z\T s_\Active\bigr\|_2^2 \ge 0,
\]
with equality only when $Z\T s_\Active = 0$, i.e.\ when $s_\Active$ lies in the span
of the active normals.
\end{proof}

\subsection{The gross-exposure hinge}

The identity has a structural reading. In $\mathrm{M}_c$ the sparsity is governed by
an \emph{explicit} constraint, the gross-exposure ball; in $\mathrm{L}_\lambda$ it is
\emph{induced} by a penalty through its subgradient. The next proposition records
when the explicit one becomes vacuous, and it is the reason the long-only frontier
looks nothing like a LASSO path even though they are the same object.

\begin{proposition}[The gross-exposure hinge]
\label{prop:hinge}
For a long-only, fully invested portfolio ($w \ge 0$, $\onevec\T w = 1$) the
gross-exposure constraint is vacuous, since $\norm{w}_1 = \onevec\T w = 1$ for all
feasible $w$. The long-only mean--variance problem is therefore a
\emph{non-negative} LASSO with the equality $\onevec\T\beta = 1$, and it traces a
frontier rather than a sparsity path. Short positions admitted under a binding cap
$\norm{w}_1 \le c$ with $c > 1$ make the constraint active, and then $\mathrm{M}_c$
and $\mathrm{L}_\lambda$ coincide as in Theorem~\ref{thm:identity}.
\end{proposition}

\begin{proof}
Immediate from $\norm{w}_1 = \sum_i w_i = \onevec\T w$ when $w \ge 0$; the budget
fixes $\norm{w}_1 = 1$ so $\norm{w}_1 \le c$ cannot bind for $c \ge 1$, and for
$c < 1$ the long-only budget is infeasible.
\end{proof}

\subsection{The third axis}

The two-dimensional risk--return picture has a \emph{third} axis, and drawing it
heads off a likely misreading. The gross-exposure formulation adds
$\norm{w}_1$, a measure of leverage; the efficient set is then a \emph{surface} in
(return, risk, leverage). The frontier and the LASSO path are perpendicular cuts of
that surface. The Critical Line Algorithm fixes the leverage structure and sweeps
the return; Brodie et al.\ \cite{brodie2009} fix the return and sweep the leverage
penalty, following the fit-against-sparsity path rather than the risk--return
frontier. Their path meets the frontier at a single point, its zero-penalty end;
raising the penalty walks off the frontier toward sparser, less-levered, slightly
riskier portfolios.

The third axis collapses when the weights are pinned, which is
Proposition~\ref{prop:hinge}: long-only and fully invested freezes leverage, and one
is back to the two-dimensional frontier with the active set driven by box bounds.
The leverage axis is genuinely live only for long--short books.

\section{A taxonomy}
\label{sec:taxonomy}

With the identity in hand, the right way to see the family is not as a list of
methods from different fields but as points in a small design space. What the members
share is the common scheme---parametric QP, affine segment, homotopy, simplex-inherited
anti-cycling---and that part is fixed, and is most of the algorithm. What varies is
captured by four features. Read them as \emph{coordinates}, not as a checklist: a
choice of all four specifies a method, and every member of the family, including
combinations no field has yet had reason to name, is some such choice.
\index{taxonomy of path methods}

\emph{Imposed versus induced active sets.} Constraints impose faces; penalties induce
them through a subgradient. Theorem~\ref{thm:identity} and
Proposition~\ref{prop:hinge} show the two are interconvertible, with the
gross-exposure ball the hinge, but their computational ergonomics differ, which is
why each field built its own version.

\emph{Monotone versus non-monotone active sets.} Forward selection and forward LARS
only \emph{add} variables; the LASSO modification and the CLA both add and drop.
The CLA corresponds to the full LASSO homotopy, not to forward LARS.

\emph{The structure of $H$.} A covariance in finance, a Gram matrix in regression, a
kernel in learning. What $H$ is fixes both interpretation and cost, and---crucially---
the homotopy never needs it explicitly (Chapter~\ref{ch:operator}).

\emph{Degeneracy.} Ties and near-rank-deficiency are the recurring practical hazard;
the inherited Bland rule guarantees determinism but, as in the simplex method, does
not tame the adversarial worst case.

Most cells of that space are occupied, several of them more than once by different
fields---the LASSO and the CLA sitting at nearly the same point and recognised as one
only through Theorem~\ref{thm:identity}. Some are barely explored: kernelised problems
under hard linear constraints, traced exactly rather than on a grid, are an imposed,
non-monotone, kernel-$H$ combination that no field has had much occasion to build, yet
the scheme constructs it with no new machinery. And because the axes are shared, a
result, an algorithmic refinement, or a piece of software attached to one cell carries
to its neighbours along them.

\section{What the identity buys, and what it does not}
\label{sec:transfer}

The force of Theorem~\ref{thm:identity} is not that two programs happen to share a
minimiser at each parameter value; it is that the two solution paths are the
\emph{same curve}. What that licenses is narrower than it first appears, and the
boundary is worth drawing precisely, because it is a good lesson in what a
mathematical identity does and does not transfer.

\paragraph{Statements about the curve transfer verbatim.} That the solution is unique
at a given point, that the two homotopies visit the same active sets in the same
order, that a segment is affine, that the path has at most so many breakpoints: each
is a geometric fact about a subset of $\R^n$, and a bound on the number of LASSO
breakpoints is at once a bound on the number of turning points of the frontier.
Algorithms and code transfer for the same reason, which is why one engine traces both.

\paragraph{Statements about a fixed value of the parameter do not.} Here the
reparametrisation $\lambda \mapsto c(\lambda)$ stops being a convenience and becomes
an obstruction. The map is $c(\lambda) = \norm{\beta(\lambda)}_1$: it is computed
\emph{from the data}, so it is not a fixed relabelling of an axis but a \emph{random}
one. Holding $\lambda$ fixed while the response varies and holding $c$ fixed while
the response varies are two different experiments, and a quantity defined by averaging
over the response is attached to the experiment, not to the curve.

Degrees of freedom is the clearest case. The general results for the lasso and
generalized lasso \cite{tibshirani2012df,zou2007df} concern the Lagrange form at a
fixed penalty; they do not by themselves yield the corresponding statement for the
constrained form at a fixed bound, because the estimator whose degrees of freedom one
would be computing is a different function of the response. The same caution applies
to post-selection inference \cite{lee2016,tibshirani2016psi} and to any risk or
model-selection criterion that integrates over the response.

\emph{The identity is an identity of curves; it is not an identity of statistical
experiments.} Results of the second kind must be proved in the parametrisation where
they are wanted. Chapter~\ref{ch:reading} takes up what survives.

\begin{notes}
This chapter follows \cite{schmelzer2026perspective}. The unconstrained
correspondence between the LASSO path and least angle regression is classical
\cite{efron2004}; the single gross-exposure case is Brodie et al.\ \cite{brodie2009};
the constrained-LASSO path under general linear equalities and inequalities is Gaines,
Kim and Zhou \cite{gaines2018}. What Theorem~\ref{thm:identity} adds is the
identification of the bound-driven frontier tracer as Markowitz's Critical Line
Algorithm---one breakpoint sequence read in both directions---together with
Proposition~\ref{prop:reg}. Rosset and Zhu \cite{rosset2007} give the general
characterisation of piecewise-linear regularised paths, and generalise in a direction
not pursued here, to piecewise-quadratic and piecewise-linear losses such as the
Huber, hinge, and quantile losses.
\end{notes}

\begin{exercises}

\begin{exercise}
Verify Theorem~\ref{thm:identity} numerically on a small design: trace the LASSO path
by any method, then solve $\mathrm{M}_c$ as a quadratic program at each
$c = \norm{\beta(\lambda)}_1$ and compare the minimisers.
\end{exercise}

\begin{exercise}
Show that $\mu \in \range(X\T)$ is necessary and sufficient for $\mu = X\T y$ to be
solvable, and give a singular $\Sig$ and $\mu$ for which it fails.
\end{exercise}

\begin{exercise}
Prove Proposition~\ref{prop:hinge} and then explain, in one sentence each, why a
long-only frontier still has breakpoints and why they are not sparsity events in the
LASSO sense.
\end{exercise}

\begin{exercise}
Give an example of a constraint system for which $Z\T s_\Active = 0$ in
Proposition~\ref{prop:reg}, so that $c$ has a flat segment. Why is this non-generic?
\end{exercise}

\begin{exercise}
Locate each of the five instances of Section~1 in the four-coordinate taxonomy of
Section~\ref{sec:taxonomy}. Which cells are unoccupied, and can you name a problem
that would occupy one?
\end{exercise}

\begin{exercise}
The forward LARS active set is monotone. Which of the four event families of
Chapter~\ref{ch:parametric} does forward LARS populate, and which does it suppress?
What is lost by suppressing it?
\end{exercise}

\begin{exercise}
Explain, using Section~\ref{sec:transfer}, why ``the number of breakpoints of the
LASSO path is at most $K$'' transfers to the frontier but ``the expected number of
nonzero coefficients at $\lambda$ equals the degrees of freedom'' does not.
\end{exercise}

\begin{exercise}[Implementation]
Take your path tracer from Chapter~\ref{ch:parametric} and drive it twice: once with
$H = \Sig$, $d = \mu$ for a small portfolio problem, and once with $H = X\T X$,
$d = X\T y$ for a small regression. Confirm that a single implementation serves both,
and compare the second against a reference LARS implementation breakpoint by
breakpoint.
\end{exercise}

\begin{exercise}[Implementation]
Trace the leverage axis: fix a return target, add the constraint $\mu\T w = \rho$,
and sweep an $\ell_1$ penalty. Plot the resulting path in (risk, leverage) and confirm
it meets the frontier only at the zero-penalty end, as Section~3.3 claims.
\end{exercise}

\end{exercises}
